% JUMP TO LINE 60, 75
\documentclass{article}
\usepackage[letterpaper,portrait,top=0.4in, left=0.6in, right=0.6in, bottom=1in]{geometry}

\usepackage{amsmath, amsfonts, amsthm, amssymb}
\usepackage{graphicx, float}
\usepackage{mathtools}
\usepackage{titlesec}
\usepackage{interval}
\usepackage{hyperref}
\usepackage[per-mode=symbol]{siunitx}
\usepackage{titling}
\usepackage{vwcol}
\usepackage{setspace}
\usepackage{empheq}
\usepackage{cancel}
\usepackage{esdiff}
\usepackage{multicol}
\usepackage{mdframed}
\usepackage{esdiff}
\usepackage{tikzsymbols}
\usepackage{multicol}
\usepackage{tikz}
\usepackage{varwidth}
\usepackage{parskip}
\usepackage{pgfplots}
\usepackage{mlmodern}
\pgfplotsset{compat=1.18}
\intervalconfig {
	soft open fences
}

\newcommand{\alignedintertext}[1]{%
  \noalign{%
    \vtop{\hsize=\linewidth#1\par
    \expandafter}%
    \expandafter\prevdepth\the\prevdepth
  }%
}

\newtheorem{lemma}{Lemma}

\renewcommand{\qedsymbol}{\Smiley[1.3]}
\newcommand*{\problem}[1]{\section*{Problem #1}}
\newcommand*{\aps}{\section*{AP Corner}}
\newcommand*{\deriv}[1][x]{\ensuremath{\dfrac{\mathrm{d}}{\mathrm{d}#1}}}
\newcommand*{\floor}[1]{\ensuremath{\lfloor #1\rfloor}}
\newcommand*{\lheqzero}{\ensuremath{\underset{\text{L'H}}{\overset{\left[\frac00\right]}{=}}}}
\newcommand*{\lheqinfty}{\ensuremath{\underset{\text{L'H}}{\overset{\left[\frac{\infty}{\infty}\right]}{=}}}}

\DeclareMathOperator{\DNE}{DNE}
\DeclareMathOperator{\sgn}{sgn}

\DeclareMathOperator{\arccsc}{arccsc}
\DeclareMathOperator{\arcsec}{arcsec}
\DeclareMathOperator{\arccot}{arccot}

%opening

\title{\vspace*{-40pt}AP Physics C -- Summer Notes}
\author{Jayden Li}
\date{\today}
% \allowdisplaybreaks
\postdisplaypenalty=100000

\begin{document}
\setstretch{1.25}
\fontsize{11pt}{12pt}\selectfont
\setlength{\abovedisplayskip}{\abovedisplayskip/2}
\setlength{\belowdisplayskip}{\belowdisplayskip/2}
\setlength{\parindent}{0pt}
\setlength{\parskip}{2ex plus 0.5ex minus 0.2ex}
\maketitle

\tableofcontents

\section{1D Kinematics}

\subsection{Displacement}

An object's \textbf{position} describes where it is at a given time, relative to a convenient reference frame.

An object's change in position is known as \textbf{displacement}. The displacement $\Delta x=x_f-x_0$, where $x_f$ is the final position, and $x_0$ is the initial position. Displacement has a direction as well as a magnitude.

\textbf{Distance} is the magnitude of the displacement, and is not signed. \textbf{Distance traveled} is the total length of the path traveled between two positions, even if the path involves a change in direction.

\subsection{Vectors, Scalars and Coordinate Systems}

In physics, a \textbf{vector} is any quantity with a direction and a magnitude. Vectors in one-dimensional motion can have direction specified by sign ($+,-$). A vector can be represented by an arrow whose length signifies magnitude.

A \textbf{scalar} is a quantity with only magnitude and no direction.

We must define a coordinate system to describe a vector's direction. Often, positive indicates right or up, and negative indicates left or down. However, this can be defined differently given the problem, but cannot be changed later.

\subsection{Time, Velocity and Speed}

\textbf{Time} is the interval over which change occurs. \textbf{Elapsed time} is the difference between the starting time of a motion and the ending time. The elapsed time $\Delta t=t_f-t_0$.

\textbf{Velocity} is the speed and direction in which an object moves. The \textbf{average velocity} is displacement (change in position) divided by elapsed time. Because displacement is a vector, velocity is also a vector. In SI units, velocity is measured in meters per second (\si{\meter\per\second}).
\begin{equation*}
    \overline v=\frac{\Delta x}{\Delta y}=\frac{x_f-x_0}{t_f-t_0}
\end{equation*}
($\overline v$ denotes the average velocity, indicated by the bar over $v$.)

\textbf{Instantaneous velocity} is an object's velocity an any given instant, or over an infinitesimally small time interval.

\textbf{Speed} is different from velocity in that speed is a scalar and has no direction. At any time, an object's \textbf{Instantaneous speed} is the magnitude of its instanteneous velocity.

\textbf{Average speed} is the distance traveled divided by elapsed time. Distance travaled can be greater than displacement, so average speed can be greater than average velocity.

\subsection{Acceleration}

\textbf{Acceleration} is the rate of change of an object's velocity. \textbf{Average acceleration} is given by:
\begin{equation*}
    \overline a=\frac{\Delta v}{\Delta t}=\frac{v_f-v_0}{t_f-t_0}
\end{equation*}
Since acceleration is velocity (\si{\meter\per\second}) divided by time (\si{\second}), the SI unit for acceleration is \si{\meter\per\second^2}.

Since velocity is a vector, acceleration is also a vector. Its direction is the direction of change in velocity, but not necessarily the direction of motion.

\textbf{Deceleration} is when an object's acceleration is opposite to the direction of its motion. This is when the object slows down.

Deceleration is not negative acceleration. Deceleration always reduces an object's speed, while negative acceleration is acceleration whose direction is negative in the chosen coordinate system.

\textbf{Instantaneous acceleration} is acceleration at a specific instant in time, or acceleration over an infinitesimally small interval of time.

\subsection{1D Motion Equations with Constant Acceleration} \label{subsec:1Dkinematics}

Let the initial time $t_0$ equal $0$, then the elapsed time $\Delta t$ equals the final time $t_f$. Let $t=\Delta t=t_f$.

Let $x_0$ be the initial position and let $x$ be the final position. Then $\Delta x=x-x_0$.

Let $v_0$ be the initial velocity and let $v$ be the final velocity. Then $\Delta v=v-v_0$.

Suppose that acceleration is constant. Then the instantaneous acceleration $a$ always equals the average acceleration $\overline a$ over any interval. Let $a$ denote acceleration.

\subsubsection{Final velocity from starting velocity and time} \label{eq:finalvelocitystarttime}
\begin{equation*}
    a
	=\frac{\Delta v}{\Delta t}
	=\frac{v-v_0}{t}
	\implies v-v_0=at
	\implies \boxed{v=v_0+at}
\end{equation*}

\subsubsection{Average velocity from starting and ending velocity} \label{eq:avgvelstartend}

Let $v(t)$ denote the object's velocity at a time $t$. If acceleration is constant, then:
\begin{equation*}
    v(s)=v_0+as
\end{equation*}
where $s$ is time, $a$ is the constant rate of acceleration, and $v_0$ is the starting velocity (from \ref{eq:finalvelocitystarttime}). Notice that the final velocity $v=v(t)=v_0+at$. We can calculate average velocity $\overline v$:
\begin{gather*}
    \overline v
	=\frac 1t \int_{0}^{t}v(s)\,\mathrm{d}s
	=\frac 1t \int_{0}^{t}\left( v_0+as \right)\mathrm{d}s
	=\frac 1t \left[v_0s+\frac{as^2}{2}\right]_{0}^{t}
	=\frac 1t \left( v_0t+\frac{at^2}{2} \right)
	=\frac{v_0+v_0+at}{2}
	=\frac{v_0+v(t)}{2} \\
	\implies \boxed{\overline v=\frac{v_0+v}{2}}
\end{gather*}

\subsubsection{Displacement, final position and velocity from average velocity}

(using equation from \ref{eq:avgvelstartend})
\begin{align*}
    \overline v 
	=\frac{\Delta x}{\Delta t}
	=\frac{x-x_0}{t}
	\implies \Aboxed{x&=x_0+\overline vt} \\ 
	\implies \Aboxed{x&=x_0+t \left( \frac{v_0+v}{2} \right)}
\end{align*}

\subsubsection{Final position when velocity is not constant}
Recall from \ref{eq:finalvelocitystarttime}, that an object's instantaneous velocity at time $s$ is:
\begin{equation*}
	v(s)=v_0+as
\end{equation*}
From Calculus, we can calculate an object's displacement from its instantaneous velocity:
\begin{gather*}
	\Delta x
	=x-x_0
	=\int_{0}^{t}v(s)\,\mathrm{d}s
	=\int_{0}^{t}\left( v_0+as \right)\mathrm{d}s
	=\left[v_0s+\frac{as^2}{2}\right]_{0}^{t}
	=v_0t+\frac{at^2}{2} \\
	\implies \boxed{x=x_0+v_0t+\frac{at^2}{2}}
\end{gather*}

\subsubsection{Final position from initial and final position}

(using \ref{eq:finalvelocitystarttime} and \ref{eq:avgvelstartend})
\begin{align*}
    v^2
	&=v \left( v_0+at \right) 
	=v_0(v_0+at)+vat
	=v_0^2+v_0at+vat
	=v_0^2+2at \left( \frac{v_0+v}{2} \right) \\
	&=v_0^2+2at\overline v
	=v_0^2+2at \left( \frac{x-x_0}{t} \right)
	=\boxed{v_0^2+2a(x-x_0)=v^2}
\end{align*}

\subsection{Falling Objects}

If air resistance and friction are negligible, all objects fall toward the center of Earth at a constant rate of accleration. The rate of acceleration differs but its average value is $g=9.80 \si{\meter\per\second}$. We can use the equations in \ref{subsec:1Dkinematics} with $a=-g$ if positive is defined as up, or $a=g$ if positive is defined as down.

\section{Forces}

\end{document}
